% !TeX program = pdflatex
\documentclass[aspectratio=169,xcolor={dvipsnames,table}]{beamer}

\input{../../_en_preambulo_beamer}
% Comandos y macros estadísticas compartidas para las presentaciones Beamer en inglés (Metropolis)

% Conjuntos numéricos básicos
\newcommand{\R}{\mathbb{R}}
\newcommand{\N}{\mathbb{N}}
\newcommand{\Q}{\mathbb{Q}}
\newcommand{\Z}{\mathbb{Z}}
\newcommand{\set}[1]{\left\{ #1 \right\}}
\newcommand{\abs}[1]{\left|#1\right|}
\newcommand{\norm}[1]{\left\| #1 \right\|}

% Comandos estadísticos y probabilísticos del curso
\newcommand{\Var}{\operatorname{Var}}
\newcommand{\cov}{\operatorname{Cov}}
\newcommand{\corr}{\rho}
\newcommand{\comb}[2]{\begin{pmatrix} #1 \\ #2 \end{pmatrix}}
\newcommand{\s}{\sigma}
\newcommand{\particion}{\mathfrak{P}}
\newcommand{\card}[1]{\#\left( #1 \right)}
\newcommand{\seq}[3]{\set{#1}_{#2}^{#3}}
\newcommand{\E}{\mathbb{E}}
\newcommand{\Prob}{\mathbb{P}}

% Entornos pedagógicos y cajas en inglés académico natural (soportando tanto nombres en inglés como en español)
\newenvironment{discussion}{%
  \begin{alertblock}{Discussion Question}%
}{%
  \end{alertblock}%
}
\newenvironment{discusion}{%
  \begin{alertblock}{Discussion Question}%
}{%
  \end{alertblock}%
}

\newenvironment{reflection}{%
  \begin{alertblock}{Classroom Reflection}%
}{%
  \end{alertblock}%
}
\newenvironment{reflexion}{%
  \begin{alertblock}{Classroom Reflection}%
}{%
  \end{alertblock}%
}

\newenvironment{paradox}{%
  \begin{alertblock}{Exploratory Paradox}%
}{%
  \end{alertblock}%
}
\newenvironment{paradoja}{%
  \begin{alertblock}{Exploratory Paradox}%
}{%
  \end{alertblock}%
}

\newenvironment{motivation}{%
  \begin{exampleblock}{Motivation: Data Science in Practice}%
}{%
  \end{exampleblock}%
}
\newenvironment{motivacion}{%
  \begin{exampleblock}{Motivation: Data Science in Practice}%
}{%
  \end{exampleblock}%
}

\newenvironment{quickchallenge}{%
  \begin{block}{Quick Team Challenge}%
}{%
  \end{block}%
}
\newenvironment{retoclase}{%
  \begin{block}{Quick Team Challenge}%
}{%
  \end{block}%
}


\title{Poisson Distribution}
\subtitle{Section 03.07 --- Arrival Processes, Law of Rare Events and Additivity}
\author[J. Castillo Colmenares]{Juliho Castillo Colmenares}
\institute[Tec de Monterrey]{Tecnologico de Monterrey}
\date{\vspace{-1.5cm}}

\begin{document}

% SLIDE 01: Title Page
\begin{frame}[plain]
	\titlepage
\end{frame}

% SLIDE 02: Roadmap
\begin{frame}{Roadmap --- Chapter 03: Discrete Random Variables}
	\vspace{-0.15cm}
	\begin{block}{Curricular Structure of Chapter 03}
		\footnotesize
		\begin{enumerate}\itemsep=0.03cm
			\item \textcolor{gray}{Section 03.01: Discrete Probability Mass Functions (PMF and Support)}
			\item \textcolor{gray}{Section 03.02: Cumulative Distribution Functions (Discrete CDF)}
			\item \textcolor{gray}{Section 03.03: Mathematical Expectation, Variance, and Moments}
			\item \textcolor{gray}{Section 03.04: Bernoulli and Binomial Distributions}
			\item \textcolor{gray}{Section 03.05: Geometric and Negative Binomial Distributions}
			\item \textcolor{gray}{Section 03.06: Hypergeometric Distribution}
			\item \textbf{\textcolor{TecRojo}{Section 03.07: Poisson Distribution}} $\leftarrow$ \emph{Current focus}
		\end{enumerate}
	\end{block}
	\vspace{-0.2cm}
	\begin{alertblock}{Learning Objective}
		\scriptsize
		Master the modeling of rare event counts in fixed time/space intervals, derive the Law of Rare Events as the bridge from Binomial to Poisson, and exploit the reproductive and conditional Binomial properties of independent Poisson random variables.
	\end{alertblock}
\end{frame}

% SLIDE 03: Motivation
\begin{frame}{From Discrete Binomial Trials to Continuous Poisson Processes}
	\vspace{-0.15cm}
	\begin{columns}[T]
		\begin{column}{0.48\textwidth}
			\begin{block}{Binomial Model: Fixed Number of Trials}
				\footnotesize
				\begin{itemize}\itemsep=0.05cm
					\item Fixed number of trials $n$ with constant success probability $p$.
					\item Discrete sample space with $n+1$ outcomes.
					\item \textbf{Limitation:} Cannot model events that occur continuously over time or space.
				\end{itemize}
			\end{block}
		\end{column}
		\begin{column}{0.48\textwidth}
			\begin{alertblock}{Poisson Model: Random Event Counts}
				\footnotesize
				\begin{itemize}\itemsep=0.05cm
					\item Random number of events in a fixed time/space window.
					\item Rate parameter $\lambda$ governs both mean and variance.
					\item \textbf{Equidispersion:} $\E[X] = \Var(X) = \lambda$.
				\end{itemize}
			\end{alertblock}
		\end{column}
	\end{columns}
\end{frame}

% SLIDE 04: Formal Definition of the PMF
\begin{frame}{Formal Definition of the Poisson Distribution}
	\vspace{-0.15cm}
	\begin{block}{Probability Mass Function (PMF)}
		\scriptsize
		A discrete random variable $X$ follows a Poisson distribution with rate $\lambda > 0$ if
		$$P(X=k)=\dfrac{e^{-\lambda}\lambda^k}{k!},\quad k\in\set{0,1,2,\dots}.$$
	\end{block}
	\vspace{-0.1cm}
	\begin{alertblock}{Interpretation and Python Parameterization}
		\scriptsize
		The factor $e^{-\lambda}$ normalizes the series, while $\lambda^k/k!$ determines its tail. In SciPy use \texttt{poisson.pmf(k, mu=lambda)} and \texttt{poisson.cdf(k, mu=lambda)}.
	\end{alertblock}
\end{frame}

% SLIDE 05: Equidispersion
\begin{frame}{Equidispersion: $\mu = \sigma^{2} = \lambda$}
	\vspace{-0.15cm}
	\begin{columns}[T]
		\begin{column}{0.48\textwidth}
			\begin{block}{Expected Value}
				\scriptsize
				$$\E[X]=\sum_{k=0}^{\infty}k\,\dfrac{e^{-\lambda}\lambda^k}{k!}=\lambda.$$
				This follows by differentiating $e^\lambda=\sum_{k=0}^{\infty}\lambda^k/k!$ and multiplying by $\lambda$.
			\end{block}
		\end{column}
		\begin{column}{0.48\textwidth}
			\begin{alertblock}{Variance}
				\scriptsize
				$$\Var(X)=\E[X^2]-(\E[X])^2=(\lambda^2+\lambda)-\lambda^2=\lambda.$$
				Thus the mean and variance coincide exactly.
			\end{alertblock}
		\end{column}
	\end{columns}
\end{frame}

% SLIDE 04B: Support and Probability Generating Function
\begin{frame}{Countable Support and the Probability Generating Function}
	\vspace{-0.15cm}
	\begin{block}{Countable Infinite Support and Unit Sum}
		\scriptsize
		The sum of Poisson probabilities over the entire support is not trivial: it requires the Taylor series identity of the exponential function.
		$$\mathcal{S}_X = \set{k \in \mathbb{Z}_{\ge 0}}, \quad \sum_{k=0}^{\infty} P(X = k) = e^{-\lambda} \sum_{k=0}^{\infty} \dfrac{\lambda^{k}}{k!} = e^{-\lambda} \cdot e^{\lambda} = 1$$
	\end{block}
	\vspace{-0.15cm}
	\begin{alertblock}{Probability Generating Function (PGF)}
		\scriptsize
		$G_{X}(t) = \E[t^{X}] = \sum_{k=0}^{\infty} t^{k} P(X=k) = e^{\lambda(t-1)}$ encodes all distributional information and will be key to proving additivity.
	\end{alertblock}
\end{frame}

% SLIDE 04C: Moment Generating Function
\begin{frame}{MGF and Characterization of the Poisson Distribution}
	\vspace{-0.15cm}
	\begin{block}{Poisson MGF}
		\scriptsize
		The Moment Generating Function $M_{X}(t) = \E[e^{tX}]$ for a Poisson variable takes a particularly simple closed form:
		$$M_{X}(t) = \sum_{k=0}^{\infty} \dfrac{e^{tk} e^{-\lambda} \lambda^{k}}{k!} = e^{-\lambda} \sum_{k=0}^{\infty} \dfrac{(\lambda e^{t})^{k}}{k!} = e^{\lambda(e^{t} - 1)}$$
	\end{block}
	\vspace{-0.15cm}
	\begin{alertblock}{Deriving Moments from the MGF}
		\scriptsize
		Raw moments are obtained by differentiation: $\E[X^{r}] = M_{X}^{(r)}(0)$. For example, $\E[X] = M_{X}'(0) = \lambda e^{0} = \lambda$ and $\E[X^{2}] = \lambda^{2} + \lambda$, confirming equidispersion.
	\end{alertblock}
\end{frame}

% SLIDE 05: Law of Rare Events
\begin{frame}{Law of Rare Events: Binomial $\to$ Poisson Limit}
	\vspace{-0.15cm}
	\begin{block}{Asymptotic Convergence Theorem}
		\scriptsize
		As $n \to \infty$ with $p \to 0$ such that $np = \lambda$ stays constant:
		$$\lim_{n \to \infty} \binom{n}{k} p^{k} (1-p)^{n-k} = \dfrac{e^{-\lambda} \lambda^{k}}{k!}$$
	\end{block}
	\vspace{-0.1cm}
	\begin{columns}[T]
		\begin{column}{0.48\textwidth}
			\begin{exampleblock}{Operational Rule}
				\scriptsize
				The approximation is excellent when $n \ge 50$, $p \le 0.05$, and $np \le 5$. Beyond these thresholds, the Total Variation Distance exceeds $10^{-3}$.
			\end{exampleblock}
		\end{column}
		\begin{column}{0.48\textwidth}
			\begin{alertblock}{Practical Use in Engineering}
				\scriptsize
				When modeling rare adverse events (clinical trials, defect counts, traffic accidents), use the Binomial for exact computation and Poisson for analytical tractability.
			\end{alertblock}
		\end{column}
	\end{columns}
\end{frame}

% SLIDE 06: Additivity
\begin{frame}{Additivity: Sum of Independent Poisson Variables}
	\vspace{-0.15cm}
	\begin{columns}[T]
		\begin{column}{0.48\textwidth}
			\begin{block}{Convolution Proof}
				\tiny
				If $X\sim\text{Pois}(\lambda_1)$ and $Y\sim\text{Pois}(\lambda_2)$ are independent, then
				\begin{align*}
				P(X+Y=m)&=\sum_{k=0}^{m}P(X=k)P(Y=m-k)\\[-0.15em]
				&=\dfrac{e^{-(\lambda_1+\lambda_2)}(\lambda_1+\lambda_2)^m}{m!}.
				\end{align*}
			\end{block}
		\end{column}
		\begin{column}{0.48\textwidth}
			\begin{alertblock}{Reproductive Property}
				\scriptsize
				The sum of independent Poisson variables is Poisson with rate $\lambda_1+\lambda_2$. This aggregates event streams without changing their distributional form.
			\end{alertblock}
		\end{column}
	\end{columns}
\end{frame}

% SLIDE 06B: Conditional Binomial
\begin{frame}{Conditional Property: Binomial Distribution under a Fixed Sum}
	\vspace{-0.15cm}
	\begin{block}{Conditional Binomial Distribution}
		\scriptsize
		If $X\sim\text{Pois}(\lambda_1)$ and $Y\sim\text{Pois}(\lambda_2)$ are independent, then
		$$P(X=k\mid X+Y=n)=\binom{n}{k}
		\left(\dfrac{\lambda_1}{\lambda_1+\lambda_2}\right)^k
		\left(\dfrac{\lambda_2}{\lambda_1+\lambda_2}\right)^{n-k}.$$
	\end{block}
	\begin{alertblock}{Queueing Interpretation}
		\scriptsize
		Given $n$ total customers served by two Poisson servers, the number served by server 1 is binomial with success probability $\lambda_1/(\lambda_1+\lambda_2)$.
	\end{alertblock}
\end{frame}

% SLIDE 07: Applications
\begin{frame}{Industrial and Data Science Applications}
	\vspace{-0.1cm}
	\begin{itemize}\itemsep=0.1cm
		\item \textbf{Queueing Theory and Telecommunications:} Call arrivals to a call center, HTTP requests to a web server, vehicle arrivals at a toll booth. The rate $\lambda$ quantifies traffic intensity. \pause
		\item \textbf{Statistical Quality Control:} Defect counts per production unit. The equidispersion property enables $c$-charts with control limits $\lambda \pm 3\sqrt{\lambda}$. \pause
		\item \textbf{Biostatistics and Epidemiology:} Genetic mutations in a DNA strand, rare disease cases in a population, or rare adverse events in clinical trials.
	\end{itemize}
\end{frame}

% SLIDE 08: Python Lab Block 1
\begin{frame}[fragile]{Python Lab: PMF and Equidispersion}
	\vspace{-0.05cm}
	\scriptsize
	Computational verification in SciPy (\texttt{03.07\_poisson\_distribution.py}) of support, normalization, and the equidispersion property:
	\vspace{0.05cm}
	{\lstset{basicstyle=\fontsize{5.5pt}{6.5pt}\selectfont\ttfamily}
	\lstinputlisting[firstline=15, lastline=37]{../../code/03_variables_aleatorias_discretas/03.07_poisson_distribution.py}}
\end{frame}

% SLIDE 09: Python Lab Block 2
\begin{frame}[fragile]{Python Lab: Law of Rare Events}
	\vspace{-0.05cm}
	\scriptsize
	$250,000$ trials simulation verifying Binomial $\to$ Poisson convergence via Total Variation Distance (TVD):
	\vspace{0.02cm}
	{\lstset{basicstyle=\fontsize{5.5pt}{6.5pt}\selectfont\ttfamily}
	\lstinputlisting[firstline=40, lastline=62]{../../code/03_variables_aleatorias_discretas/03.07_poisson_distribution.py}}
\end{frame}

% SLIDE 10: Python Lab Block 3
\begin{frame}[fragile]{Python Lab: Additivity and Conditional Binomial}
	\vspace{-0.05cm}
	\scriptsize
	Numerical verification by simulation of Poisson additivity and the conditional Binomial distribution:
	\vspace{0.02cm}
	{\lstset{basicstyle=\fontsize{5pt}{6pt}\selectfont\ttfamily}
	\lstinputlisting[firstline=72, lastline=92]{../../code/03_variables_aleatorias_discretas/03.07_poisson_distribution.py}}
\end{frame}

% SLIDE 11: Python Lab Terminal Output
\begin{frame}[fragile]{Python Lab: Terminal Output}
	\vspace{-0.1cm}
	\begin{block}{}
		\fontsize{4.5pt}{5.2pt}\selectfont
		\begin{verbatim}
=== Block 1: Poisson PMF & Equidispersion ===
Lambda = 3.0 (calls/min) | mu = sigma^2 = lambda
PMF total probability: 0.999999329614
Equidispersion property verified: True
P(X=0)=0.049787 | P(X=1)=0.149361 | P(X=2)=0.224042

=== Block 2: Law of Rare Events (Binomial -> Poisson) ===
  n=    50 p=0.0400 | TVD=0.009270
  n=   500 p=0.0040 | TVD=0.000905
  n=  2000 p=0.0010 | TVD=0.000226
  n= 10000 p=0.0002 | TVD=0.000045
Monte Carlo (N=250,000): mean=2.0018 | var=1.9971

=== Block 3: Additivity & Conditional Binomial ===
Z = X + Y: emp mean=7.9992 | emp var=7.9888
Conditional X | (X+Y=8) with p_cond = 0.3750:
  k=0: 0.0222 | 0.0233 | k=3: 0.2856 | 0.2816
  k=5: 0.1019 | 0.1014 | k=8: 0.0003 | 0.0004
		\end{verbatim}
	\end{block}
\end{frame}

% SLIDE 12: Exercise Level 1 (Statement)
\begin{frame}{Closing Section and Modular Perspective}
	\vspace{-0.15cm}
	\begin{block}{Synthesis on the Law of Rare Events}
		\scriptsize
		The Poisson distribution quantifies both uncertainty and expected value with a single parameter $\lambda$. This equidispersion property is the foundation of statistical process control and queueing modeling.
	\end{block}
	\vspace{0.02cm}
	\pause
	\begin{alertblock}{Key Lessons and Next Step}
		\begin{itemize}\itemsep=0.04cm
			\item \textbf{Equidispersion:} $\E[X] = \Var(X) = \lambda$ for prediction and variability monitoring.
			\item \textbf{Additivity:} Sum of independent Poisson is Poisson; ideal for aggregating event flows.
			\item \textbf{Conditional Binomial:} Under fixed total, the distribution of a component is Binomial.
			\item \textbf{Next Session:} \textcolor{TecRojo}{Multinomial Distribution}.
		\end{itemize}
	\end{alertblock}
\end{frame}

% SLIDE 18: Modular Perspective
\begin{frame}{Modular Perspective --- The Next Challenge}
	\begin{columns}[T]
		\begin{column}{0.48\textwidth}
			\begin{block}{Beyond Fixed Event Counts}
				\scriptsize
				Poisson counts events in a time or space interval at a fixed rate $\lambda$. The next challenge is to model waiting times between events rather than counts.
			\end{block}
		\end{column}
		\begin{column}{0.48\textwidth}
			\begin{alertblock}{Section 03.08}
				\scriptsize
				We will study the Multinomial Distribution and its applications to several-category experiments.
			\end{alertblock}
		\end{column}
	\end{columns}
\end{frame}

\end{document}
